% Section 6: Channel 3 -- Market Making
\label{sec:channel3}

\subsection{Setup}

Channel 3 studies how AI signal correlation among market makers reduces the effective diversification of liquidity supply. Consider $N$ market makers, each using an AI-calibrated model to post bid-ask quotes in a single asset with value $V \sim \mathcal{N}(\bar{v}, \sigma_V^2)$. In the \citet{avellaneda2008} stochastic control framework, each market maker $i$ posts a reservation price
\begin{equation}\label{eq:reservation-price}
  r_i(t) = S_i(t) - q_i\gamma_i\sigma_i^2(T-t),
\end{equation}
where $S_i(t)$ is the midpoint estimate, $q_i$ is inventory, $\gamma_i$ is risk aversion, and $\sigma_i^2$ is the estimated asset variance. The reservation price~\eqref{eq:reservation-price} summarises how each market maker's quoting behaviour depends on her model-derived estimate of asset variance. When market makers share AI-calibrated models, their estimates inherit the $\rho$-parameterised correlation structure:
\begin{equation}\label{eq:correlated-calibrations}
  \sigma_i^2 = \sigma_V^2 + \varepsilon_i^\sigma, \quad \varepsilon_i^\sigma = \sqrt{\rho}\,\eta_\sigma + \sqrt{1-\rho}\,\xi_i^\sigma,
\end{equation}
where $\eta_\sigma \sim \mathcal{N}(0,\sigma_{\text{cal}}^2)$ is the common AI calibration error and $\xi_i^\sigma \sim \mathcal{N}(0,\sigma_{\text{cal}}^2)$ is idiosyncratic. The calibration structure in~\eqref{eq:correlated-calibrations} implies that a negative AI model error $\eta_\sigma$ simultaneously raises all market makers' variance estimates, pushing them toward the withdrawal threshold. Similarly, midpoint estimates $S_i(t) = V + \sqrt{\rho}\,\eta_S + \sqrt{1-\rho}\,\xi_i^S$ share the same correlation structure. Market maker $i$ withdraws from the market (sets spread to infinity) when estimated variance exceeds a tolerance threshold $\sigma_{\max}^2$.

The market-making framework descends from the \citet{kyle1985} inventory-management tradition, adapted by \citet{avellaneda2008}, rather than the \citet{glosten1985} adverse-selection framework. The relevant distinction is that our market makers manage inventory risk from correlated position limits, not adverse selection from informed order flow. The present channel applies this framework to the case where $N$ market makers share a common AI calibration, producing correlated spread-widening functions. The key question is: how many independent liquidity providers does a market with $N$ AI-calibrated market makers effectively have?

\subsection{The N-Effective Liquidity Fragility Index}

Each market maker's withdrawal decision is a Bernoulli random variable $W_i$ with mean $p_w$ and variance $p_w(1-p_w)$. The correlation of withdrawal decisions inherits the correlation of calibration errors: $\mathrm{Corr}(W_i, W_j) = \rho$ for $i \neq j$. The variance of the average withdrawal rate $\bar{W} = (1/N)\sum_{i=1}^N W_i$ is
\begin{equation}\label{eq:withdrawal-variance}
  \mathrm{Var}(\bar{W}) = \frac{p_w(1-p_w)(1+(N-1)\rho)}{N}.
\end{equation}
The factor $(1+(N-1)\rho)$ in~\eqref{eq:withdrawal-variance} captures the variance inflation from correlated withdrawal decisions. When $\rho = 0$, the variance is $p_w(1-p_w)/N$, as in the independent case. When $\rho > 0$, the variance rises because a common AI calibration error pushes all market makers toward the same exit decision. For $N_{\text{eff}}$ independent market makers, $\mathrm{Var}(\bar{W}) = p_w(1-p_w)/N_{\text{eff}}$. Equating these expressions yields

\begin{equation}\label{eq:neff}
  N_{\text{eff}}(\rho) = \frac{N}{1+(N-1)\rho}.
\end{equation}

\begin{proposition}[Liquidity Fragility Index]\label{prop:neff}
  The effective number of independent liquidity providers $N_{\text{eff}}(\rho)$ defined in~\eqref{eq:neff} satisfies:

  \partlabel{(i)} $N_{\text{eff}}(0) = N$ (full independence) and $N_{\text{eff}}(1) = 1$ (single effective market maker).

  \partlabel{(ii)} $N_{\text{eff}}$ is strictly decreasing in $\rho$:
  \begin{equation}\label{eq:neff-decreasing}
    \frac{dN_{\text{eff}}}{d\rho} = -\frac{N(N-1)}{(1+(N-1)\rho)^2} < 0.
  \end{equation}
  The monotonicity in~\eqref{eq:neff-decreasing} confirms that any increase in AI calibration similarity reduces the effective number of independent liquidity providers. The rate of decline is largest at $\rho = 0$, where the denominator equals one, and attenuates as $\rho$ increases.

  \partlabel{(iii)} $N_{\text{eff}}$ is strictly convex in $\rho$:
  \begin{equation}\label{eq:neff-convex}
    \frac{d^2N_{\text{eff}}}{d\rho^2} = \frac{2N(N-1)^2}{(1+(N-1)\rho)^3} > 0.
  \end{equation}

  \partlabel{(iv)} $N_{\text{eff}}(1/(N-1)) = N/2$: half the market makers are effectively lost at $\rho = 1/(N-1)$.

  \prooflabel{Proof.} Direct computation from~\eqref{eq:neff}. Properties (ii) and (iii) follow by differentiation; (iv) follows by substitution. $\square$
\end{proposition}

The convexity established in~\eqref{eq:neff-convex} is economically significant. The marginal loss of independence is greatest at low correlation, where the diversification benefit is still present to be lost. Consider a market with $N = 100$ market makers and $\rho = 0.1$. It has $N_{\text{eff}} \approx 9.2$: only nine percent of the nominal diversity is effective. A small increase in AI model similarity, from $\rho = 0.05$ to $\rho = 0.10$, halves the effective number of independent liquidity providers. The nominal count is misleading. The effective count is what determines fragility.

While direct evidence on AI calibration correlation among market makers is still emerging, the VaR model homogeneity documented by \citet{danielsson2012}, where common risk constraints cause market participants to act in harmony during stress, provides a direct precedent for the mechanism modelled here. \citet{kirilenko2017} document the fragility of algorithmic intermediation during the May 2010 flash crash. More recently, \citet{liu2025} show empirically that strategic homogeneity among algorithmic traders amplifies market volatility and increases mini-crash risk, and the \citet{fsb2024} identifies correlated AI model use as a key financial stability vulnerability.

The $N_{\text{eff}}(\rho)$ formula extends the \citet{greenwood2011} stock fragility measure from the demand side to the supply side. Greenwood and Thesmar define stock fragility as the covariance of holders' demand shocks, capturing how correlated mutual fund flows create instability. The present measure captures the analogous supply-side phenomenon: correlated AI calibrations create instability in liquidity supply. The mechanism is algorithmic similarity rather than ownership overlap.

\subsection{Equilibrium Spread and the No-Equilibrium Threshold}

In the competitive market-making equilibrium, the equilibrium half-spread is proportional to the inverse of the effective number of liquidity providers:
\begin{equation}\label{eq:spread}
  s^*(\rho) = \frac{s_0}{N_{\text{eff}}(\rho)} = s_0 \cdot \frac{1+(N-1)\rho}{N},
\end{equation}
where $s_0$ is the single-market-maker (monopoly) half-spread. The spread is strictly increasing in $\rho$ with $s^*(0) = s_0/N$ (competitive spread) and $s^*(1) = s_0$ (monopoly spread). Let $Q$ denote the total order flow per period and $\kappa_{\text{inv}}$ the inventory cost coefficient, which scales the variance cost of holding correlated positions.

\begin{proposition}[No-Equilibrium Threshold]\label{prop:rho-star-star}
  Assume $N > 1$, $\sigma_V^2 > 0$, $\gamma > 0$, outside option $\pi_0 > 0$, and the parameter condition $Qs_0/N^2 > \gamma\sigma_V^2\kappa_{\text{inv}}$ (positive marginal revenue net of inventory costs). Then there exists a threshold $\rho^{**} \in (0,1)$ such that:

  \partlabel{(i)} For $\rho < \rho^{**}$: an interior market-making equilibrium with finite spreads exists.

  \partlabel{(ii)} For $\rho > \rho^{**}$: no interior market-making equilibrium with finite spreads exists. All market makers withdraw simultaneously, and the market enters the Pagano (1989) low-liquidity trap.

  The threshold $\rho^{**}$ is the solution to the participation constraint at equality: $\mathrm{Revenue}(\rho^{**}) = \mathrm{Cost}(\rho^{**}) + \pi_0$, where in the linear specification
  \begin{equation}\label{eq:rho-star-star}
    \rho^{**} = \frac{1}{N-1}\left(\frac{\pi_0}{Q s_0/N^2 - \gamma\sigma_V^2\kappa_{\text{inv}}} - 1\right),
  \end{equation}
  \prooflabel{Proof sketch.} Market maker $i$ participates when expected profit $E[\pi_i^{MM}(\rho)] \geq \pi_0$. Revenue per market maker is $Qs^*(\rho)/N$, where $s^*$ is the equilibrium spread from~\eqref{eq:spread}. Inventory cost depends on the covariance of market-maker positions: correlated AI calibrations produce correlated inventories, creating concentrated systemic risk. In the linear specification of~\eqref{eq:rho-star-star}, the participation constraint reduces to an affine condition in $(1+(N-1)\rho)$. The threshold $\rho^{**}$ is the unique solution where net profit (revenue minus inventory cost minus outside option $\pi_0$) equals zero. The parameter condition ensures the denominator is positive and $\rho^{**} \in (0,1)$. When this condition fails, $\rho^{**} = 0$ and no equilibrium exists for any positive $\rho$. For $\rho > \rho^{**}$, the outside option dominates, and rational market makers exit simultaneously. $\square$
\end{proposition}

The threshold $\rho^{**}$ is analogous to the \citet{pagano1989} low-liquidity trap, where the expectation of thin markets deters participation, validating the expectation. Here, the expectation of correlated market-maker withdrawal raises effective adverse selection costs for remaining market makers, reinforcing the withdrawal. Above $\rho^{**}$, the anticipated cost of correlated inventory risk exceeds the revenue from market making, and no market maker can profitably remain. The simultaneous exit of all identically positioned market makers validates the high-volatility estimate that triggered the exit, completing a self-fulfilling liquidity withdrawal.

\subsection{Nesting the Danielsson-Shin-Zigrand Framework}

\begin{proposition}[DSZ Limiting Case]\label{prop:dsz}
  The \citet{danielsson2012} common VaR framework is the limiting case $\rho = 1$ of the Channel 3 model. Specifically:

  \partlabel{(i)} At $\rho = 1$: $N_{\text{eff}}(1) = 1$, and all $N$ market makers act as a single representative market maker with the DSZ common VaR constraint.

  \partlabel{(ii)} The DSZ procyclical spiral (VaR breach $\to$ deleveraging $\to$ price impact $\to$ higher measured volatility $\to$ further VaR breaches) corresponds to the feedback between $s^*(1) = s_0$ and correlated inventory cost at $\rho = 1$.

  \partlabel{(iii)} The present model generalises DSZ by allowing $\rho \in [0,1)$, with $N_{\text{eff}}(\rho)$ as a continuous measure of proximity to the DSZ limiting case.

  \prooflabel{Proof.} At $\rho = 1$, all calibration errors are identical: $\varepsilon_i^\sigma = \eta_\sigma$ for all $i$. All market makers have the same reservation price and withdrawal threshold, acting identically. The system reduces to a single market maker. $\square$
\end{proposition}

The DSZ framework identified that common regulatory constraints (VaR) create endogenous systemic risk by forcing simultaneous deleveraging. The Channel 3 model extends this insight. AI model homogeneity creates the same common constraint endogenously, through shared calibrations rather than regulation. The $N_{\text{eff}}$ formula provides a continuous measure of proximity to the DSZ limiting case, making the qualitative insight of \citet{danielsson2012} quantitatively tractable.

The contrast with \citet{cespa2025} is instructive. Cespa and Vives show that market opacity creates strategic complementarity in liquidity demand, yielding multiple equilibria where liquidity evaporates suddenly. Their mechanism operates through information opacity: a market maker quotes narrower spreads when she believes others will also provide liquidity. The present mechanism operates through algorithmic similarity: market makers withdraw simultaneously because the common AI calibration error drives them to the same threshold at the same time. The two mechanisms could compound in practice, but they are formally distinct.

Channel 3 in isolation understates the danger. When $N_{\text{eff}}$ falls, the effective signal correlation in the coordination game rises, pushing the economy closer to the Channel 1 uniqueness boundary. A higher crisis probability destroys the option value of private information, reducing $\mu_I^*$ and further concentrating market-maker reliance on the common AI signal. Each channel feeds the others. Section~\ref{sec:amplification} formalises this positive feedback loop and derives the compound fragility threshold that no single-channel analysis can detect.
% --- Clarity pass complete: 2026-03-10 ---
% --- Flow pass complete: 2026-03-11 ---
% --- Technical audit complete: 2026-03-11 ---
